
\chapter{Objetivos}

\section{Objetivo general}

Desarrollar y evaluar un sistema de control para el robot cuadrúpedo Nova Spot
Micro que combine una marcha nominal con una política de aprendizaje por
refuerzo como capa correctiva acotada, con el propósito de mejorar la
estabilidad y la repetibilidad durante la postura y la locomoción tipo paso y
gateo en una superficie plana y bajo condiciones controladas de laboratorio.


\section{Objetivos específicos}


\begin{itemize}

\item Caracterizar la plataforma Nova Spot Micro y desarrollar los modelos
cinemático y dinámico necesarios para su representación en simulación y para la
generación de referencias de movimiento.

\item Diseñar e implementar una arquitectura electrónica y de software que
permita accionar y supervisar el robot de forma segura, respetando los límites
eléctricos, mecánicos y articulares de la plataforma.

\item Implementar una estrategia convencional de locomoción basada en una
máquina de estados finitos y referencias nominales para postura, paso y gateo.

\item Entrenar en simulación e integrar progresivamente en el robot una política
de aprendizaje por refuerzo que genere correcciones pequeñas y acotadas sobre
las referencias nominales, sin comandar directamente las señales PWM ni
sustituir el supervisor de seguridad.

\item Comparar experimentalmente el desempeño de la marcha nominal con y sin la
capa correctiva aprendida mediante métricas de estabilidad, seguimiento y
repetibilidad, tanto en simulación como en el robot físico.

\end{itemize}
